% The collision cone in three dimensions: a right circular cone touching
% the sphere of radius R = r_A + r_B along a circle.
% Sphere centre (5,1.2), radius 1.1: d = 5.142, axis 13.5 deg, theta = 12.35 deg;
% tangent circle at distance 4.907 along the axis, radius 1.0745.
\begin{tikzpicture}[scale=1.0]
\coordinate (O) at (0,0);
\coordinate (C) at (5,1.2);
% cone surface (near part) as a light fill to the far cross-section
\fill[sbOrange!12] (O) -- (6.633,3.213) -- (7.369,0.149) -- cycle;
\draw[sbOrange,thick] (O) -- (4.520,2.190);
\draw[sbOrange,thick] (O) -- (5.022,0.100);
\draw[sbOrange,thick,dashed] (4.520,2.190) -- (6.633,3.213);
\draw[sbOrange,thick,dashed] (5.022,0.100) -- (7.369,0.149);
\draw[sbOrange!70,thin,rotate around={13.5:(7.001,1.681)}] (7.001,1.681) ellipse [x radius=0.47,y radius=1.576];
% the sphere and its equator
\draw[sbregion,fill=sbOrange!25] (C) circle (1.1);
\draw[sbOrange!80!black,dashed,thin] (C) ellipse [x radius=1.1,y radius=0.32];
\node[font=\small] at ($(C)+(0.2,-0.55)$) {$B$};
% tangent circle
\draw[sbOrange!80!black,thick,rotate around={13.5:(4.771,1.145)}] (4.771,1.145) ellipse [x radius=0.32,y radius=1.0745];
% axis, radius, angle
\draw[dashed,sbGray] (O) -- (7.6,1.824);
\draw[sbGray] (C) -- ++(60:1.1) node[pos=1.3,sbannot] {$R$};
\draw[sbGray] (1.5,0) arc[start angle=0,end angle=25.85,radius=1.5];
\node[sbannot] at (1.85,0.42) {$2\theta$};
\node[sbannot] at (2.6,1.05) {$d$};
% a dangerous relative velocity
\draw[sbvec,draw=sbRed] (O) -- (22:3.3) node[pos=1,above,sbannot,text=sbRed] {$\vel$};
\node[sbagentA] at (O) {$A$};
\node[sbannot] at (3.55,2.75) {tangent circle};
\draw[sbGray,thin] (3.9,2.6) -- (4.45,2.1);
\node[sbannot,text=sbOrange!80!black] at (6.9,-0.35) {$CC_{A|B}$};
\end{tikzpicture}
